\section{本章小结}
\label{sec:quant-summary}

本章开头问：每个Laplace读数只存2个比特时，4个代表值放在哪里、读数怎样归属、平均误差多大。答案依次是：从均匀刻度出发，平均失真为$0.514$；按两条必要条件交替改进，得到阈值$\pm1.594$、重构值$\pm0.594$与$\pm2.594$，平均失真降到$0.352$；若允许变长记录，同一量化器的索引平均只需$1.805$ bit，而改用均匀刻度配合熵编码，在相近的平均长度下平均失真可以降到$0.21$左右。

\paragraph{三样东西决定三件事}
一个量化器由区域与重构值组成，它们共同决定平均失真$D(Q)$；索引的编码方式决定平均描述长度；是否加入共享的随机扰动，决定误差与读数之间的关系。三者分工不同：无损的索引编码不改变重构值，从而不改变失真；减性抖动不减小平均失真，只改变误差的结构。

\paragraph{平方误差下的两条条件}
固定区域时，最好的重构值是区域的质心，这来自平行轴恒等式~\eqref{eq:quant-parallel-axis}，也说明此时量化输出是给定索引下的条件期望；固定重构值时，最好的区域是最近邻区域，标量时阈值取中点，向量时得到Voronoi区域。最优量化器必须满足这两条，但满足它们的量化器未必最优。Lloyd--Max算法及其向量推广交替执行两步，平均失真单调不增，可能停在局部最优，结果依赖初值。

\paragraph{构造之间的比较}
本章的各种构造分别改变量化的不同环节，作用与代价各不相同。

\begin{BookInlineTable}
 \centering
 \begin{tabular}{lll}
  \toprule
  构造 & 得到什么 & 代价\\
  \midrule
  均匀量化 & 编码最简单 & 不适应分布，需饱和\\
  Lloyd--Max量化 & 按分布调整区间 & 需知道分布，局部最优\\
  向量量化 & 利用相关与形状 & 码本规模$2^{dR}$\\
  索引的熵编码 & 平均长度近$H(J)$ & 变长记录，需码表\\
  熵约束量化 & 按平均码长设计 & 目标含码长项\\
  减性抖动 & 误差独立于读数 & 共享随机数\\
  \bottomrule
 \end{tabular}
 \caption{本章各种构造的作用与代价。}
 \label{tab:quant-summary}
\end{BookInlineTable}

表~\ref{tab:quant-summary}汇总了各种构造。比较它们时，必须固定同一个源分布、同一个失真函数和同一种描述资源：重构点数、定长索引长度与平均索引长度是三种不同的预算，在一种预算下最好的方案，在另一种预算下未必最好。

\paragraph{误差的结构}
有限范围的均匀量化把误差分成两部分：范围内的舍入误差不超过半个步长，高分辨率下约为$\Delta^{2}/12$乘以范围内的概率；范围外的过载误差无界，其大小与尾部概率由读数的尾部决定。步长和范围的选择是这两部分之间的权衡，点数增加时范围也必须随之变宽。确定性量化的误差是读数的函数，“独立均匀噪声”只是高分辨率下的近似；只有在减性抖动下，它才精确成立。

\paragraph{有限方案与渐近极限}
本章的所有方案都是有限的：给定维数和点数，可以执行，也可以算出代价。它们不回答一个更根本的问题：允许把任意长的读数块放在一起描述时，给定每读数的比特数，平均失真最低能到多少。第~\ref{chap:rate-distortion}章用重构条件分布上的优化刻画这一极限，并把本章的有限量化器放到同一个标准下，衡量它们与极限之间的差距。
